\subsection{Бизнес-процесс в модели TO-BE}\label{subsec:business-process-in-to-be-model}

В отличие от первой версии бизнес-процесса, операции по доставке изменённых веб-сервисов в целевые среды выполняет подсистема автоматизации сборки и развёртывания.\ При этом первичный анализ рисков, постановка задачи, проверка исходного кода и согласование реализованного функционала остаются в зоне ответственности каждого из участников процесса.

Отправка изменений во временную ветку автоматически запускает сборку изменённых веб-сервисов и развёртывание отдельного экземпляра приложения в тестовой среде.\ Если в ходе сборки или развёртывания возникают ошибки, разработчик получает информацию о них и вносит необходимые исправления, после чего повторная отправка изменений запускает эти операции заново.

Проверка исходного кода на соответствие техническому заданию в таком варианте бизнес-процесса существенно упрощается, поскольку работающий экземпляр уже собран и позволяет наглядно протестировать изменения. Кроме того, автоматическое развёртывание по мере слияний в основную ветку позволяет добиться постоянной синхронизации между актуальным исходным кодом и пре-продакшеном.

После выполнения всех требований заказчика руководитель команды в соответствии с установленным расписанием релизов самостоятельно инициирует развёртывание согласованной версии приложения в продакшен через интерфейс подсистемы автоматизации.

Модель бизнес-процесса после внедрения автоматизации представлена на рисунке~\ref{fig:business-process-to-be}. Состав поставщиков, входов, основных процессов, выходов и потребителей модели TO-BE представлен в таблице~\ref{tab:sipoc-to-be} по методологии SIPOC.

\begin{landscape}
	\begin{table}[p]
		\centering
		\caption{SIPOC-модель бизнес-процесса TO-BE}\label{tab:sipoc-to-be}
		\footnotesize
		\setstretch{1}
		\setlength{\tabcolsep}{5pt}
		\renewcommand{\arraystretch}{1.15}
		\begin{tabularx}{\linewidth}{|Y{0.9}|Y{0.95}|Y{1.3}|Y{1.05}|Y{0.8}|}
			\hline
			\textbf{Поставщики} & \textbf{Входы} & \textbf{Процессы} & \textbf{Выходы} & \textbf{Потребители} \\
			\hline
			Заказчик; менеджер проекта &
			Обратная связь: функциональная проблема или дефект &
			1. Первичный анализ рисков, ограничений и возможности реализации руководителем команды &
			Решение о согласовании изменения или отказе от его реализации &
			Руководитель команды \\
			\hline
			Руководитель команды &
			Согласованное изменение; выявленные ограничения &
			2. Подготовка технического задания, создание задачи в трекере и назначение разработчиков &
			Техническое задание; задача с назначенными исполнителями &
			Разработчик \\
			\hline
			Руководитель команды; разработчик &
			Техническое задание; исходный код компонентов системы &
			3. Внесение изменений и отправка кода во временную ветку репозитория &
			Изменённый код; событие отправки изменений в ветку &
			Подсистема автоматизации \\
			\hline
			Разработчик &
			Код временной ветки; событие её обновления &
			4. Автоматическая сборка и развёртывание тестового экземпляра при обновлении ветки &
			Тестовый экземпляр либо сообщение об ошибках сборки и развёртывания &
			Руководитель команды; разработчик \\
			\hline
			Руководитель команды; разработчик; подсистема автоматизации &
			Техническое задание; исходный код; работающий тестовый экземпляр &
			5. Проверка кода и тестирование руководителем; слияние изменений в основную ветку &
			Проверенный код в основной ветке либо замечания для доработки &
			Подсистема автоматизации; разработчик \\
			\hline
			Руководитель команды; заказчик &
			Код основной ветки; событие слияния; требования заказчика &
			6. Автоматическая сборка и развёртывание в пре-продакшене; проверка и согласование заказчиком &
			Экземпляр в пре-продакшене; согласование изменений либо замечания &
			Заказчик; руководитель команды; разработчик \\
			\hline
			Заказчик; руководитель команды &
			Согласованная версия приложения; команда на выпуск по расписанию релизов &
			7. Запуск выпуска руководителем через интерфейс подсистемы; автоматическое развёртывание в продакшене &
			Обновлённая система в рабочей среде &
			Заказчик; пользователи системы \\
			\hline
		\end{tabularx}
	\end{table}
\end{landscape}
